\section*{Ethical Considerations}

Our findings reveal that reasoning traces in language models, while often seen as a step toward transparency or interpretability, can introduce vulnerabilities with potential safety and privacy implications. We show that these traces are difficult to steer in a controlled way, can contain unsafe content, and are relatively easy to extract, even in unintended scenarios. These characteristics raise concerns about their possible misuse, such as inferring sensitive information or manipulating model behavior for malicious purposes.

At the same time, we view this work as a contribution to the responsible development of language technologies. By systematically analyzing and exposing these issues, our goal is to raise awareness within the research and practitioner communities. Understanding the limitations and risks of reasoning traces is an important step toward developing models that are safer, more reliable, and more aligned with user expectations.

There is a clear dual-use aspect to this work. While it may draw attention to specific weaknesses, it also enables researchers, developers, and users to better understand and anticipate the kinds of failures and threats that may arise. We have aimed to present these findings in a way that supports transparency and encourages mitigation efforts, rather than facilitating direct misuse.

\section*{Acknowledgments}
This work was supported by the NAVER corporation.
